\providecommand{\rank}{\operatorname{r}}
\providecommand{\RR}{\mathbb R}
\providecommand{\CC}{\mathbb C}
\providecommand{\KK}{\mathbb K}

% ---------- Exercise 29 ----------
\setcounter{enumi}{28}
\item {\bfseries\boldmath
Let
}

\[
A=\begin{bmatrix}3&2&0\\0&0&1\\0&1&3\end{bmatrix},\quad B=\begin{bmatrix}1&2&0\\0&-2&1\\0&1&1\end{bmatrix},\quad C=\begin{bmatrix}1&0\\0&2\\-1&0\end{bmatrix}.
\]

\begin{enumerate}[label=\alph*)]

\item Compute the inverse of $B$.

\begin{answer}
Here
\[
 A=\begin{bmatrix}3&2&0\\0&0&1\\0&1&3\end{bmatrix},\quad
 B=\begin{bmatrix}1&2&0\\0&-2&1\\0&1&1\end{bmatrix},\quad
 C=\begin{bmatrix}1&0\\0&2\\-1&0\end{bmatrix}.
\]

On $[B\mid I_3]$, use
\[
 F_2\gets-\tfrac12F_2,\quad F_3\gets F_3-F_2,\quad
 F_3\gets\tfrac23F_3,\quad F_2\gets F_2+\tfrac12F_3,\quad
 F_1\gets F_1-2F_2.
\]
Then $B$ is invertible and
\[
 B^{-1}=\begin{bmatrix}1&2/3&-2/3\\0&-1/3&1/3\\0&1/3&2/3\end{bmatrix}.
\]
\end{answer}

\item Solve $AX=2X+BC+C$ for $X\in M_{3\times2}(\RR)$.

\begin{answer}
Since $A-2I_3=B$, the equation is equivalent to
$BX=BC+C$. Multiplying on the left by $B^{-1}$ gives the unique solution
\[
 X=C+B^{-1}C
 =\begin{bmatrix}1&0\\0&2\\-1&0\end{bmatrix}
 +\begin{bmatrix}5/3&4/3\\-1/3&-2/3\\-2/3&2/3\end{bmatrix}
 =\begin{bmatrix}8/3&4/3\\-1/3&4/3\\-5/3&2/3\end{bmatrix}.
\]
Substitution gives $BX=BC+C$, hence the original equation.
\end{answer}
\end{enumerate}
